\section{Skenario Pengujian Fungsional (\textit{Black Box Testing})}
Pengujian dilakukan dengan metode \textit{Black Box} yang berfokus pada fungsionalitas input-output tanpa melihat kode internal. Berikut adalah rekapitulasi hasil pengujian fitur utama:

\begin{table}[H]
\centering
\caption{Hasil Pengujian Fungsional Sistem}
\label{tab:functional_test}
\begin{tabular}{|p{0.5cm}|p{4cm}|p{5cm}|p{2.5cm}|}
    \hline
    \textbf{No} & \textbf{Skenario Pengujian} & \textbf{Hasil yang Diharapkan} & \textbf{Status} \\ \hline
    1 & Pendaftaran Agen Baru & Data tersimpan di database, password terenkripsi (bcrypt), dan redirect ke login. & \textbf{Berhasil} \\ \hline
    2 & Validasi Input Sewa & Sistem menolak jika tanggal \textit{check-out} $<$ tanggal \textit{check-in}. & \textbf{Berhasil} \\ \hline
    3 & Pencatatan Log History & Setiap perubahan pada kolom komentar atau waktu terekam di tabel \texttt{rental\_logs}. & \textbf{Berhasil} \\ \hline
    4 & Ekspor Data & Sistem menghasilkan file CSV yang berisi rekapitulasi data penyewa dengan format yang benar. & \textbf{Berhasil} \\ \hline
    5 & Deteksi Anomali (Simulasi) & Sistem memberikan tanda (\textit{flag}) pada data dummy yang memiliki pola edit berulang kali. & \textbf{Berhasil} \\ \hline
    \end{tabular}
\end{table}